\newpage

\section{Заключение}

В работе рассмотрена последовательная линия развития алгоритма Frank--Wolfe для задач машинного обучения и крупномасштабной выпуклой оптимизации.
Сначала были описаны модификации FW, использующие предыдущие направления в задаче равновесного распределения транспортных потоков.
Затем эти идеи были перенесены на задачи машинного обучения, сводимые к гладкой выпуклой оптимизации на компактном множестве.
После этого была рассмотрена стохастическая блочно-координатная версия, уменьшающая стоимость линейного оракула за счет обновления только части блоков.
Наконец, была сформулирована децентрализованная блочно-координатная схема, объединяющая линейные оракулы, консенсус и отслеживание градиента.

Основные результаты работы:
\begin{enumerate}
    \item Описаны CFW, BFW, FFW, WFFW и NFW как модификации FW, ускоряющие сходимость за счет памяти о предыдущих направлениях.
    \item Показано, что CFW можно применять в машинном обучении без явного вычисления гессиана, используя приближение через разности градиентов.
    \item Приведена оценка сходимости CFW по зазору Frank--Wolfe, совпадающая по порядку со стандартной оценкой FW.
    \item Описан SOFW/BCFW как блочно-координатный метод для транспортных сетей, уменьшающий стоимость итерации за счет стохастического выбора источников.
    \item Сформулирован DBCFW.
    Это децентрализованное расширение блочного FW с градиентным отслеживанием.
\end{enumerate}

\subsection{Итоги по главам}

В первой главе было показано, что основная слабость классического FW в транспортной
задаче связана не с линейным оракулом, а с геометрией направления.
CFW, BFW и NFW используют прошлые направления для уменьшения зигзагообразного поведения,
а FFW и WFFW решают похожую задачу через сглаживание.
Экспериментальные результаты показывают, что память о нескольких направлениях
может существенно ускорять уменьшение относительного зазора.

Во второй главе эта идея была перенесена в машинное обучение.
Ключевой технический шаг состоит в отказе от явного гессиана.
Вместо произведения $\nabla^2 f(x_k)d_{k-1}$ используется разность соседних градиентов.
Благодаря этому CFW и NFW становятся применимыми к логистической регрессии с ограничением
на $l_1$-норму.
Эксперименты на Mushrooms и MNIST показывают, что выигрыш особенно заметен, когда
безусловное решение находится далеко от границы допустимого множества.

В третьей главе основное внимание было перенесено с направления на стоимость итерации.
SOFW обновляет только часть источников транспортной сети и тем самым уменьшает стоимость
линейного оракула.
Теоретически метод соответствует блочно-координатному FW и сохраняет стандартный порядок
сходимости в ожидании.
Практически же он выигрывает на крупных сетях, где полный запуск кратчайших путей становится
главной вычислительной нагрузкой.

В четвертой главе блочная идея была соединена с децентрализованной оптимизацией.
DBCFW сохраняет механизм консенсуса и градиентного отслеживания из DFW, но заменяет полный
линейный оракул блочным.
Если выбрать все блоки, метод возвращается к обычному DFW.
Если убрать децентрализацию, получается блочный FW.
Поэтому DBCFW является естественным завершающим звеном рассматриваемой хронологии.

\subsection{Ограничения и дальнейшая работа}

Ограничение рассмотренной линии методов состоит в том, что каждая модификация улучшает
свой вычислительный аспект: направление шага, стоимость линейного оракула или распределенную
архитектуру.
Единый метод, который одновременно использует сопряженные направления, стохастический выбор
блоков и децентрализованный консенсус, остается отдельной исследовательской задачей.
Перспективным направлением является объединение стохастического выбора блоков с сопряженными
направлениями NFW, а также экспериментальная проверка DBCFW на распределенных задачах
машинного обучения.
